\documentclass[a4paper,12pt]{article}

\usepackage[T2A]{fontenc}
\usepackage[utf8]{inputenc}
\usepackage[russian]{babel}
\usepackage{amsmath,amssymb,amsthm,mathtools}
\usepackage{helvet}
\renewcommand{\sfdefault}{cmss}
\renewcommand{\familydefault}{\sfdefault}
\usepackage{icomma}
\usepackage[a4paper,margin=20mm,headheight=25pt,footskip=14mm]{geometry}
\usepackage[nopatch=footnote]{microtype}
\usepackage{tikz}
\usetikzlibrary{calc,arrows.meta,positioning,shapes.geometric}
\usepackage{tkz-euclide}
\usepackage{tabularx,booktabs}
\usepackage{enumitem}
\usepackage{hyperref}
\usepackage{parskip}
\usepackage{fancyhdr}
\usepackage{setspace}
\usepackage{xcolor}

\definecolor{accent}{HTML}{008080}
\definecolor{darkaccent}{HTML}{004D4D}
\definecolor{textgray}{HTML}{333333}
\definecolor{softgray}{HTML}{6B7280}
\definecolor{linegray}{HTML}{D6DEDE}

\hypersetup{
    colorlinks=true,
    linkcolor=darkaccent,
    urlcolor=darkaccent
}

\color{textgray}
\onehalfspacing
\setlength{\parindent}{0pt}

\setlist[itemize]{
    leftmargin=1.6em,
    itemsep=0.25em,
    topsep=0.35em
}
\setlist[enumerate]{
    leftmargin=1.8em,
    itemsep=0.45em,
    topsep=0.4em
}

\pagestyle{fancy}
\fancyhf{}
\fancyhead[L]{\small\textcolor{softgray}{
    Геометрия: трапеции, окружности, прямоугольные треугольники
}}
\fancyhead[R]{\small\textcolor{softgray}{София Кальней}}
\fancyfoot[C]{\small\thepage}
\renewcommand{\headrulewidth}{0.3pt}
\renewcommand{\headrule}{%
    \hbox to\headwidth{%
        \color{linegray}\leaders\hrule height \headrulewidth\hfill
    }%
}
\renewcommand{\footrulewidth}{0pt}

\newcommand{\Section}[1]{%
    \vspace{0.4em}%
    {\Large\bfseries\color{darkaccent}#1}\par
    \vspace{0.25em}
    \noindent
    \textcolor{linegray}{\rule{\linewidth}{0.6pt}}
    \vspace{0.5em}%
}

\newcommand{\Subsection}[1]{%
    \vspace{0.3em}%
    {\large\bfseries\color{accent}#1}\par
    \vspace{0.25em}%
}

\newcommand{\Key}[1]{%
    \textbf{\textcolor{darkaccent}{#1}}%
}

% -------------------------------------------------
% Стили блок-схем
% -------------------------------------------------

\tikzset{
    flowstart/.style={
        ellipse,
        draw=darkaccent,
        line width=0.9pt,
        minimum width=42mm,
        minimum height=11mm,
        align=center,
        font=\small
    },
    flowaction/.style={
        rectangle,
        rounded corners=2mm,
        draw=accent,
        line width=0.9pt,
        text width=84mm,
        minimum height=14mm,
        align=center,
        font=\small,
        inner sep=3mm
    },
    flowdecision/.style={
        diamond,
        aspect=2.2,
        draw=darkaccent,
        line width=0.9pt,
        text width=48mm,
        align=center,
        font=\small,
        inner sep=1.5mm
    },
    flowarrow/.style={
        -{Stealth[length=3mm,width=2mm]},
        line width=0.9pt,
        draw=darkaccent
    },
    flowlabel/.style={
        font=\small,
        fill=white,
        inner sep=1pt,
        text=darkaccent
    }
}

% -------------------------------------------------
% Специальные страницы для крупных блок-схем
% -------------------------------------------------

\newenvironment{flowchartpagelandscape}{%
    \clearpage
    \pdfpagewidth=420mm
    \pdfpageheight=297mm
    \paperwidth=420mm
    \paperheight=297mm
    \newgeometry{
        layoutwidth=420mm,
        layoutheight=297mm,
        left=18mm,
        right=18mm,
        top=15mm,
        bottom=15mm,
        headheight=12pt,
        headsep=0pt,
        footskip=8mm
    }
    \thispagestyle{empty}
}{%
    \clearpage
    \restoregeometry
    \pdfpagewidth=210mm
    \pdfpageheight=297mm
    \paperwidth=210mm
    \paperheight=297mm
}

\newenvironment{flowchartpageportrait}{%
    \clearpage
    \pdfpagewidth=297mm
    \pdfpageheight=420mm
    \paperwidth=297mm
    \paperheight=420mm
    \newgeometry{
        layoutwidth=297mm,
        layoutheight=420mm,
        left=20mm,
        right=20mm,
        top=18mm,
        bottom=18mm,
        headheight=12pt,
        headsep=0pt,
        footskip=8mm
    }
    \thispagestyle{empty}
}{%
    \clearpage
    \restoregeometry
    \pdfpagewidth=210mm
    \pdfpageheight=297mm
    \paperwidth=210mm
    \paperheight=297mm
}

\begin{document}

% =================================================
% ТИТУЛЬНЫЙ ЛИСТ
% =================================================

\thispagestyle{empty}

\begin{center}

\vspace*{23mm}

{\small\textcolor{softgray}{ОГЭ по математике · геометрия}}\par

\vspace{10mm}

{\Huge\bfseries\color{darkaccent}Чек-лист}\par

\vspace{4mm}

{\Large\bfseries
Трапеции, окружности и прямоугольные треугольники
}\par

\vspace{5mm}

{\large
Касательные · биссектрисы · высота к гипотенузе · описанная окружность
}\par

\vspace{18mm}

{\large\bfseries София Кальней}\par

{\normalsize программа: ОГЭ}\par

\vfill

{\normalsize \today}\par

\vspace{4mm}

{\small
Лёвин Артём Александрович эксклюзивно для Софии Кальней
}\par

\vspace{15mm}

\begin{tikzpicture}[x=1cm,y=1cm]
    \draw[
        accent,
        line width=1.1pt
    ]
    (0,0)
    .. controls (3.0,0.85) and (7.2,-0.65) ..
    (10.5,0.08)
    .. controls (12.6,0.52) and (14.8,0.30) ..
    (16.2,0.02);
\end{tikzpicture}

\end{center}

\clearpage

% =================================================
\Section{1. Главная идея занятия}

Практически вся работа строится вокруг одного принципа:
\Key{старайтесь переводить конфигурацию к треугольникам,
особенно к прямоугольным}.
Окружность и трапеция дают полезные дополнительные построения,
после которых начинают работать стандартные теоремы.

\begin{itemize}

    \item Если есть \Key{касательная к окружности},
    проведите радиус в точку касания.

    \item Если окружность \Key{вписана в угол},
    соедините вершину угла с центром окружности:
    получится биссектриса.

    \item Если окружность \Key{вписана в многоугольник},
    отмечайте равные касательные, выходящие из одной вершины.

    \item В равнобедренной трапеции проводите высоту
    к большему основанию:
    она сразу даёт прямоугольный треугольник
    и удобные отрезки на основании.

    \item Если трапеция вписана в окружность,
    полезно провести диагональ и перейти к треугольнику,
    вписанному в ту же окружность.

\end{itemize}

% =================================================
\Section{2. Что нужно знать наизусть}

\Subsection{2.1. Касательные и вписанная окружность}

\begin{enumerate}

    \item Радиус, проведённый в точку касания,
    перпендикулярен касательной:
    \[
        OT\perp l.
    \]

    \item Касательные, проведённые из одной точки
    к окружности, равны.

    Если из точки $A$ проведены касательные
    $AB$ и $AC$, то
    \[
        AB=AC.
    \]

    \item Центр окружности, вписанной в угол,
    лежит на биссектрисе этого угла.

    \item Для четырёхугольника,
    в который вписана окружность,
    суммы длин противоположных сторон равны:
    \[
        AB+CD=BC+AD.
    \]

    Это следствие равенства касательных,
    выходящих из каждой вершины.

\end{enumerate}

% -------------------------------------------------
\Subsection{2.2. Равнобедренная трапеция}

Пусть $ABCD$ --- равнобедренная трапеция,
\[
AD\parallel BC,
\qquad
AD>BC.
\]

Из вершины $C$ опущена высота
\[
CH\perp AD.
\]

Тогда
\[
AH=\frac{AD+BC}{2},
\qquad
DH=\frac{AD-BC}{2}.
\]

Поскольку средняя линия трапеции равна
\[
m=\frac{AD+BC}{2},
\]
получаем важное соотношение
\[
\boxed{AH=m}.
\]

Если боковая сторона равна $a$,
а угол при большем основании равен $30^\circ$,
то в прямоугольном треугольнике $CDH$
\[
CH=a\sin30^\circ=\frac a2,
\]
\[
DH=a\cos30^\circ=\frac{a\sqrt3}{2}.
\]

% -------------------------------------------------
\Subsection{2.3. Прямоугольный треугольник: высота к гипотенузе}

Пусть в прямоугольном треугольнике $ABC$
угол $B$ прямой,
а $BH$ --- высота к гипотенузе $AC$.

Обозначим
\[
AH=p,
\qquad
HC=q,
\qquad
AC=c,
\qquad
BH=h.
\]

Тогда
\[
\boxed{h^2=pq},
\qquad
\boxed{AB^2=cp},
\qquad
\boxed{BC^2=cq}.
\]

Также из равенства площадей
\[
\frac12 AB\cdot BC
=
\frac12 AC\cdot BH
\]
следует
\[
\boxed{
h=\frac{AB\cdot BC}{AC}
}.
\]

% -------------------------------------------------
\Subsection{2.4. Описанная окружность}

Для треугольника со стороной $a$,
лежащей напротив угла $\alpha$,
верна расширенная теорема синусов:
\[
\frac{a}{\sin\alpha}=2R.
\]

Отсюда
\[
\boxed{
R=\frac{a}{2\sin\alpha}
}.
\]

Особенно полезен частный случай:
\[
\alpha=30^\circ.
\]

Так как
\[
\sin30^\circ=\frac12,
\]
получаем
\[
\boxed{R=a},
\]
где $a$ --- сторона,
лежащая напротив угла $30^\circ$.

Если сторона треугольника является диаметром
описанной окружности,
то угол, опирающийся на этот диаметр,
равен $90^\circ$.

\clearpage

% =================================================
\Section{3. Наглядная схема: высота к гипотенузе}

\begin{center}

\begin{tikzpicture}[scale=1.05]

    \tkzDefPoint(0,0){A}
    \tkzDefPoint(8,0){C}
    \tkzDefPoint(3,4){B}

    \tkzDefPointBy[projection=onto A--C](B)
    \tkzGetPoint{H}

    \tkzDrawPolygon[
        line width=1pt,
        color=accent
    ](A,B,C)

    \tkzDrawSegment[
        line width=1pt,
        color=darkaccent
    ](B,H)

    \tkzMarkRightAngle[size=0.28](A,B,C)
    \tkzMarkRightAngle[size=0.28](B,H,C)

    \tkzDrawPoints(A,B,C,H)

    \tkzLabelPoints[below](A,H,C)
    \tkzLabelPoints[above](B)

    \node[
        below=8mm of H,
        text=darkaccent,
        font=\small
    ]{
        $BH^2=AH\cdot HC$
    };

\end{tikzpicture}

\end{center}

\textbf{Почему формула работает.}
Треугольники $ABH$ и $BCH$ подобны по двум углам.

Из отношения сходственных сторон
\[
\frac{AH}{BH}
=
\frac{BH}{HC}
\]
получаем
\[
BH^2=AH\cdot HC.
\]

Это одна из ключевых формул занятия:
она позволяет находить высоту
или один из отрезков гипотенузы
без введения дополнительных неизвестных.

% =================================================
\Section{4. Типовая модель 1: прямоугольная трапеция с вписанной окружностью}

Пусть в прямоугольную трапецию $ABCD$
вписана окружность радиуса $r$,
\[
AB\parallel CD,
\qquad
AD\perp AB,
\]
а меньшее основание
\[
AB=\frac43r.
\]

Найдём все стороны трапеции.

\begin{center}

\begin{tikzpicture}[scale=1.45]

    \tkzDefPoint(0,2){A}
    \tkzDefPoint(1.3333,2){B}
    \tkzDefPoint(4,0){C}
    \tkzDefPoint(0,0){D}

    \tkzDefPoint(1,1){O}

    \tkzDefPoint(1,2){P}
    \tkzDefPoint(1.6,1.8){T}
    \tkzDefPoint(1,0){Q}
    \tkzDefPoint(0,1){S}

    \tkzDrawPolygon[
        line width=1pt,
        color=accent
    ](A,B,C,D)

    \tkzDrawCircle[
        line width=0.9pt,
        color=darkaccent
    ](O,P)

    \tkzDrawSegments[
        line width=0.85pt,
        color=darkaccent
    ](
        O,P
        O,T
        O,Q
        O,S
    )

    \tkzMarkRightAngle[size=0.20](A,D,C)
    \tkzMarkRightAngle[size=0.18](O,P,B)
    \tkzMarkRightAngle[size=0.18](O,T,C)

    \tkzDrawPoints(A,B,C,D,O,P,T,Q,S)

    \tkzLabelPoints[above left](A)
    \tkzLabelPoint[above right](B){$B$}
    \tkzLabelPoint[above left](P){$P$}
    \tkzLabelPoints[right](C,T)
    \tkzLabelPoints[below](D,Q)
    \tkzLabelPoints[left](S)
    \tkzLabelPoints[below right](O)

\end{tikzpicture}

\end{center}

\begin{enumerate}

    \item Окружность касается обоих параллельных оснований.
    Поэтому расстояние между основаниями
    равно диаметру:
    \[
        AD=2r.
    \]

    \item Из вершины $A$ выходят две равные касательные.
    Поэтому отрезки от $A$
    до точек касания на $AB$ и $AD$
    равны $r$.

    Следовательно,
    \[
        BP
        =
        AB-AP
        =
        \frac43r-r
        =
        \frac13r.
    \]

    \item Из вершины $B$ касательные равны:
    \[
        BT=BP=\frac13r.
    \]

    \item Центр $O$ лежит на биссектрисах
    углов $B$ и $C$.

    Соседние углы трапеции в сумме дают
    \[
        \angle ABC+\angle BCD=180^\circ.
    \]

    Поэтому
    \[
        \angle OBC+\angle BCO
        =
        \frac12\left(
            \angle ABC+\angle BCD
        \right)
        =
        90^\circ.
    \]

    Значит,
    \[
        \angle BOC=90^\circ.
    \]

    Треугольник $BOC$ прямоугольный,
    а $OT$ --- высота к гипотенузе $BC$.

    \item По свойству высоты,
    проведённой к гипотенузе,
    \[
        OT^2=BT\cdot TC.
    \]

    Подставляем
    \[
        OT=r,
        \qquad
        BT=\frac13r:
    \]
    \[
        r^2
        =
        \frac13r\cdot TC.
    \]

    Отсюда
    \[
        TC=3r.
    \]

    \item Наклонная боковая сторона:
    \[
        BC
        =
        BT+TC
        =
        \frac13r+3r
        =
        \frac{10}{3}r.
    \]

    \item Из вершины $C$ выходят равные касательные:
    \[
        CQ=CT=3r.
    \]

    Из вершины $D$ также выходят равные касательные:
    \[
        DQ=DS=r.
    \]

    Поэтому
    \[
        CD
        =
        CQ+QD
        =
        3r+r
        =
        4r.
    \]

\end{enumerate}

Итак,
\[
\boxed{
AB=\frac43r,
\qquad
BC=\frac{10}{3}r,
\qquad
CD=4r,
\qquad
AD=2r
}.
\]

\textbf{Проверка.}
Для четырёхугольника с вписанной окружностью
\[
AB+CD=BC+AD.
\]

Действительно,
\[
\frac43r+4r
=
\frac{10}{3}r+2r
=
\frac{16}{3}r.
\]

\clearpage

% =================================================
\Section{5. Типовая модель 2: равнобедренная трапеция в окружности}

Пусть боковая сторона равнобедренной трапеции
равна $a$,
средняя линия равна $b$,
а углы при большем основании
равны $30^\circ$.

Требуется найти радиус окружности,
описанной около трапеции.

Проведём диагональ $AC$
и высоту
\[
CH\perp AD.
\]

Тогда треугольник $ACD$
вписан в ту же окружность,
что и трапеция.

\begin{center}

\begin{tikzpicture}[scale=1.05]

    \tkzDefPoint(0,0){A}
    \tkzDefPoint(5.464,0){D}
    \tkzDefPoint(1.732,1){B}
    \tkzDefPoint(3.732,1){C}
    \tkzDefPoint(3.732,0){H}

    \tkzDrawPolygon[
        line width=1pt,
        color=accent
    ](A,B,C,D)

    \tkzDrawSegment[
        line width=0.9pt,
        color=darkaccent
    ](A,C)

    \tkzDrawSegment[
        line width=0.9pt,
        color=darkaccent
    ](C,H)

    \tkzMarkRightAngle[size=0.22](C,H,A)

    \tkzMarkAngle[size=0.52](C,D,A)
    \tkzLabelAngle[pos=0.80](C,D,A){$30^\circ$}

    \tkzDrawPoints(A,B,C,D,H)

    \tkzLabelPoints[below](A,H,D)
    \tkzLabelPoints[above](B,C)

\end{tikzpicture}

\end{center}

Так как $CH$ --- высота равнобедренной трапеции,
\[
AH
=
\frac{AD+BC}{2}
=
b.
\]

В прямоугольном треугольнике $CDH$
\[
CH
=
a\sin30^\circ
=
\frac a2.
\]

Следовательно,
из прямоугольного треугольника $ACH$
\[
AC^2
=
AH^2+CH^2
=
b^2+\frac{a^2}{4}.
\]

Поэтому
\[
AC
=
\sqrt{
    b^2+\frac{a^2}{4}
}.
\]

В треугольнике $ACD$
сторона $AC$ лежит напротив угла
\[
\angle ADC=30^\circ.
\]

По расширенной теореме синусов
\[
R
=
\frac{AC}{2\sin30^\circ}
=
AC.
\]

Итак,
\[
\boxed{
R=
\sqrt{
    b^2+\frac{a^2}{4}
}
}.
\]

% -------------------------------------------------
\Subsection{Полезное наблюдение}

Если известны $a$ и $b$,
можно восстановить оба основания трапеции.

В прямоугольном треугольнике $CDH$
\[
DH
=
a\cos30^\circ
=
\frac{a\sqrt3}{2}.
\]

Для равнобедренной трапеции
\[
DH
=
\frac{AD-BC}{2}.
\]

Следовательно,
\[
AD-BC=a\sqrt3.
\]

Средняя линия даёт
\[
AD+BC=2b.
\]

Решая систему,
получаем
\[
\boxed{
AD=b+\frac{a\sqrt3}{2}
},
\]
\[
\boxed{
BC=b-\frac{a\sqrt3}{2}
}.
\]

Чтобы меньшее основание было положительным,
необходимо
\[
\boxed{
b>\frac{a\sqrt3}{2}
}.
\]

\clearpage

% =================================================
% БЛОК-СХЕМА
% =================================================

\begin{flowchartpageportrait}

\begin{center}

{\Large\bfseries\color{darkaccent}
Алгоритм: трапеция и окружность
}\par

\vspace{10mm}

\begin{tikzpicture}[
    node distance=5mm and 15mm
]

    \node[flowstart] (start) {
        Прочитайте условие и сделайте точный чертёж
    };

    \node[
        flowdecision,
        below=of start
    ] (circle) {
        Есть окружность и касания?
    };

    \node[
        flowaction,
        below left=8mm and 15mm of circle
    ] (tangent) {
        Проведите радиусы в точки касания;
        отметьте перпендикуляры;
        выпишите равные касательные;
        проведите биссектрисы к центру
    };

    \node[
        flowaction,
        below right=8mm and 15mm of circle
    ] (trap) {
        Проверьте тип трапеции.
        Для равнобедренной проведите
        высоту к большему основанию
    };

    \node[
        flowaction,
        below=32mm of circle
    ] (tri) {
        Найдите прямоугольный треугольник
        или создайте его дополнительным построением:
        высотой, диагональю, радиусом
    };

    \node[
        flowdecision,
        below=of tri
    ] (alt) {
        Есть высота к гипотенузе?
    };

    \node[
        flowaction,
        below left=8mm and 15mm of alt
    ] (metric) {
        Используйте
        $h^2=pq$,
        $a^2=cp$,
        $b^2=cq$
        или
        $h=ab/c$
    };

    \node[
        flowaction,
        below right=8mm and 15mm of alt
    ] (pyth) {
        Используйте Пифагора,
        тригонометрию угла $30^\circ$
        и свойства трапеции
    };

    \node[
        flowdecision,
        below=30mm of alt
    ] (Rneed) {
        Нужно найти радиус
        описанной окружности?
    };

    \node[
        flowaction,
        below left=8mm and 15mm of Rneed
    ] (sine) {
        Перейдите к вписанному треугольнику
        и примените
        \[
            R=\frac{a}{2\sin\alpha}
        \]
    };

    \node[
        flowaction,
        below right=8mm and 15mm of Rneed
    ] (collect) {
        Соберите искомые стороны
        из найденных отрезков
    };

    \node[
        flowaction,
        below=28mm of Rneed
    ] (check) {
        Проверьте длины, знаки,
        применимость формул
        и согласованность с чертежом
    };

    \node[
        flowstart,
        below=of check
    ] (finish) {
        Запишите ответ
    };

    \draw[flowarrow]
        (start) -- (circle);

    \draw[flowarrow]
        (circle)
        --
        node[
            flowlabel,
            pos=0.28,
            left
        ]{да}
        (tangent);

    \draw[flowarrow]
        (circle)
        --
        node[
            flowlabel,
            pos=0.28,
            right
        ]{нет}
        (trap);

    \draw[flowarrow]
        (tangent)
        |-
        (tri.west);

    \draw[flowarrow]
        (trap)
        |-
        (tri.east);

    \draw[flowarrow]
        (tri) -- (alt);

    \draw[flowarrow]
        (alt)
        --
        node[
            flowlabel,
            pos=0.28,
            left
        ]{да}
        (metric);

    \draw[flowarrow]
        (alt)
        --
        node[
            flowlabel,
            pos=0.28,
            right
        ]{нет}
        (pyth);

    \draw[flowarrow]
        (metric)
        |-
        (Rneed.west);

    \draw[flowarrow]
        (pyth)
        |-
        (Rneed.east);

    \draw[flowarrow]
        (Rneed)
        --
        node[
            flowlabel,
            pos=0.28,
            left
        ]{да}
        (sine);

    \draw[flowarrow]
        (Rneed)
        --
        node[
            flowlabel,
            pos=0.28,
            right
        ]{нет}
        (collect);

    \draw[flowarrow]
        (sine)
        |-
        (check.west);

    \draw[flowarrow]
        (collect)
        |-
        (check.east);

    \draw[flowarrow]
        (check) -- (finish);

\end{tikzpicture}

\end{center}

\end{flowchartpageportrait}

% =================================================
\Section{6. Типичные ошибки и контрольные вопросы}

\begin{itemize}

    \item
    \Key{Радиус к касательной.}
    Перпендикуляр возникает только
    в точке касания.

    \item
    \Key{Равные касательные.}
    Сравниваются отрезки,
    выходящие из одной и той же внешней точки.

    \item
    \Key{Биссектриса.}
    Центр вписанной окружности
    лежит на биссектрисе угла,
    образованного двумя касательными.

    \item
    \Key{Высота к гипотенузе.}
    Формула
    \[
        h^2=pq
    \]
    применима только к высоте,
    проведённой из вершины прямого угла
    к гипотенузе.

    \item
    \Key{Проекции катетов.}
    В формулах
    \[
        a^2=cp,
        \qquad
        b^2=cq
    \]
    каждому катету соответствует
    его собственная проекция на гипотенузу.

    \item
    \Key{Диаметр.}
    Если сторона треугольника
    является диаметром окружности,
    противоположный угол равен $90^\circ$.

    \item
    \Key{Радиус через угол.}
    В формуле
    \[
        R=\frac{a}{2\sin\alpha}
    \]
    сторона $a$ обязана лежать
    напротив угла $\alpha$.

    \item
    \Key{Угол $30^\circ$.}
    Если напротив $30^\circ$
    лежит сторона $a$, то
    \[
        R=a.
    \]

    \item
    \Key{Равнобедренная трапеция.}
    Формулы
    \[
        AH=\frac{AD+BC}{2},
        \qquad
        DH=\frac{AD-BC}{2}
    \]
    требуют равнобедренности трапеции.

\end{itemize}

% -------------------------------------------------
\Subsection{Быстрая самопроверка перед решением}

\begin{enumerate}

    \item Где можно провести перпендикуляр?

    \item Какие касательные равны?

    \item Какие биссектрисы проходят
    через центр окружности?

    \item Какой прямоугольный треугольник
    уже есть или может быть построен?

    \item Есть ли высота к гипотенузе?

    \item Можно ли заменить задачу о трапеции
    задачей о вписанном треугольнике?

    \item Какая сторона лежит
    напротив известного угла?

\end{enumerate}

\clearpage

% =================================================
\Section{7. Тренировочный блок — Путь А}

\textbf{Условия.}
Решайте без готовых шаблонов перед глазами.
Для каждой задачи сначала выполните чертёж
и подпишите все данные.
Решения в этом блоке не приводятся.

\begin{enumerate}[label=\textbf{\arabic*.}]

    \item
    Из точки $P$, лежащей вне окружности,
    проведены касательные $PA$ и $PB$.
    Известно, что
    \[
        PA=7.
    \]
    Найдите $PB$.

    \item
    В прямоугольную трапецию
    вписана окружность радиуса $6$.
    Меньшее основание трапеции равно $8$.
    Найдите все четыре стороны трапеции.

    \item
    Боковая сторона равнобедренной трапеции
    равна $10$,
    средняя линия равна $12$,
    а углы при большем основании
    равны $30^\circ$.
    Найдите радиус окружности,
    описанной около трапеции.

    \item
    Равнобедренная трапеция $ABCD$
    вписана в окружность,
    \[
        AD\parallel BC.
    \]
    Большее основание
    \[
        AD=26
    \]
    является диаметром окружности,
    а
    \[
        BC=14.
    \]
    Найдите высоту трапеции,
    боковую сторону
    и диагональ $AC$.

    \item
    Боковая сторона равнобедренной трапеции
    равна $8$,
    углы при большем основании
    равны $30^\circ$,
    а радиус описанной окружности
    равен $10$.
    Найдите среднюю линию
    и оба основания трапеции.

\end{enumerate}

\clearpage

% =================================================
\Section{8. Ответы к тренировочному блоку}

\renewcommand{\arraystretch}{1.35}

\begin{table}[ht]

\centering

\begin{tabularx}{\linewidth}{@{}>{\bfseries}c X@{}}

\toprule
№ & Ответ \\
\midrule

1 &
$PB=7$.
\\

2 &
Меньшее основание $8$,
большее основание $24$,
перпендикулярная боковая сторона $12$,
наклонная боковая сторона $20$.
\\

3 &
$R=13$.
\\

4 &
Высота
$2\sqrt{30}$,
боковая сторона
$2\sqrt{39}$,
диагональ
$AC=2\sqrt{130}$.
\\

5 &
Средняя линия
$2\sqrt{21}$;
основания
\[
2\sqrt{21}-4\sqrt3
\]
и
\[
2\sqrt{21}+4\sqrt3.
\]
\\

\bottomrule

\end{tabularx}

\end{table}

\vfill

\begin{center}
\textcolor{softgray}{\small
Ключевой ориентир:
окружность и трапеция должны привести Вас
к треугольнику,
с которым уже можно работать
стандартными теоремами.
}
\end{center}

\end{document}